% =============================================================================
% Hauptdatei des Scientific Writing Kit
% =============================================================================
% Nicht direkt kompilieren. Aufruf immer über /pdf bzw.
%   node .claude/kit/werkzeuge/pdf.mjs            (fertiges PDF)
%   node .claude/kit/werkzeuge/pdf.mjs entwurf    (mit Platzhaltern und Wasserzeichen)
% pdf.mjs kopiert diese Vorlage nach .lokal/build/, erzeugt dort kit-daten.tex,
% kapitel-liste.tex und die generierten Verzeichnisse und kompiliert.
% Klasse: KOMA scrreprt (Standard) oder tudscrreprt (vorlage: tudscr unter Layout in .arbeit/einstellungen.md).
% =============================================================================

\input{kit-daten}
\expandafter\documentclass\expandafter[\kitKlassenoptionen]{\kitKlasse}
% =============================================================================
% Präambel des Scientific Writing Kit
% =============================================================================
% Herstellerzone: Wird per /update erneuert. Eigene Anpassungen gehören nach
% .arbeit/latex/eigene-praeambel.tex (wird am Ende geladen, falls vorhanden).
%
% Engine: LuaLaTeX (MiKTeX, TeX Live, MacTeX). XeLaTeX (Tectonic) als Fallback.
% Alle \kit...-Makros kommen aus kit-daten.tex, das .claude/kit/werkzeuge/pdf.mjs aus
% .arbeit/einstellungen.md erzeugt. Nie von Hand ändern, sondern die Einstellungen.
% =============================================================================

\usepackage{iftex}
\ifPDFTeX
  \errmessage{Bitte mit LuaLaTeX oder XeLaTeX kompilieren (node .claude/kit/werkzeuge/pdf.mjs)}
\fi

% --- Schrift ---------------------------------------------------------------
\usepackage{fontspec}
\ifx\kitSchrift\empty
  % Standard: Latin Modern, in jeder TeX-Distribution enthalten.
\else
  \setmainfont{\kitSchrift}
\fi

% --- Sprache ---------------------------------------------------------------
\ifkitEnglisch
  \usepackage[ngerman,main=english]{babel}
\else
  \usepackage[english,main=ngerman]{babel}
\fi
\usepackage[autostyle=true]{csquotes}

% --- Seite und Absatz ------------------------------------------------------
\usepackage[
  top=\kitRandOben, bottom=\kitRandUnten,
  inner=\kitRandInnen, outer=\kitRandAussen,
  includefoot, footskip=1.2cm
]{geometry}
\usepackage{setspace}
\setstretch{\kitZeilenabstand}
\usepackage{microtype}

% --- Kopf- und Fußzeile (KOMA) ---------------------------------------------
\usepackage[automark,headsepline=false]{scrlayer-scrpage}
\clearpairofpagestyles
\cfoot*{\pagemark}
\ohead{\headmark}
\setkomafont{pagehead}{\normalfont\small\itshape}
\setkomafont{pagenumber}{\normalfont}

% --- Überschriften ---------------------------------------------------------
% Serifen statt KOMA-Serifenlos, Größen wie in der erprobten Bachelor-Vorlage.
\setkomafont{disposition}{\normalfont\bfseries}
\ifdefined\chapter\setkomafont{chapter}{\Large}\fi
\setkomafont{section}{\large}
\setkomafont{subsection}{\normalsize}
\setkomafont{subsubsection}{\normalsize}
\ifdefined\chapter
  \RedeclareSectionCommand[beforeskip=0pt, afterskip=18pt]{chapter}
\fi
\RedeclareSectionCommand[beforeskip=-18pt plus -4pt minus -2pt, afterskip=6pt]{section}
\RedeclareSectionCommand[beforeskip=-12pt plus -3pt minus -2pt, afterskip=6pt]{subsection}
\setcounter{secnumdepth}{3}
\setcounter{tocdepth}{2}

% --- Verzeichnisse ---------------------------------------------------------
% Abstände im Inhaltsverzeichnis und Schutz gegen Überlappung langer Titel mit
% der Seitenzahl (erprobt in der Vorgängerarbeit).
\makeatletter
\renewcommand{\@pnumwidth}{2.2em}
\renewcommand{\@tocrmarg}{3.5em}
\makeatother
\ifdefined\chapter
  \DeclareTOCStyleEntry[beforeskip=10pt plus 1pt]{chapter}{chapter}
\fi
\DeclareTOCStyleEntry[beforeskip=2pt]{section}{section}

% --- Mathematik, Einheiten, Chemie -----------------------------------------
\usepackage{amsmath}
\usepackage{amssymb}
\usepackage{siunitx}
\ifkitEnglisch
  \sisetup{locale=UK, per-mode=symbol, separate-uncertainty=true}
\else
  \sisetup{locale=DE, per-mode=symbol, separate-uncertainty=true}
\fi
\usepackage[version=4]{mhchem}
\usepackage{chemfig}

% --- Abbildungen und Tabellen ----------------------------------------------
\usepackage{graphicx}
% Bilder liegen im Build-Ordner (.lokal/build/) unter abbildungen/, Fallback: Projektwurzel (../../).
\graphicspath{{./}{abbildungen/}{../../}}
\usepackage{xcolor}
\usepackage{booktabs}
\usepackage{longtable}
% Pandoc ab 3.x setzt bei Tabellen ohne Beschriftung \LTcaptype{none}: Zähler muss existieren.
\makeatletter
\@ifundefined{c@none}{\newcounter{none}}{}
\makeatother
\usepackage{array}
\usepackage{tabularx}
\usepackage{calc}
\usepackage[font=small, labelfont=bf, format=hang]{caption}
\usepackage{float}
\usepackage{enumitem}
% Abbildungen und Tabellen durchgehend nummerieren (1, 2, 3 statt 3.1, 3.2).
\ifdefined\chapter
  \counterwithout{figure}{chapter}
  \counterwithout{table}{chapter}
  \counterwithout{equation}{chapter}
\fi
% Abstände zwischen Gleitobjekten und Text.
\setlength{\textfloatsep}{20pt plus 4pt minus 4pt}
\setlength{\intextsep}{20pt plus 4pt minus 4pt}
\setlength{\floatsep}{16pt plus 4pt minus 4pt}
% Nicht mehr als eine Seite nur mit Gleitobjekten erzwingen.
\renewcommand{\topfraction}{0.85}
\renewcommand{\textfraction}{0.1}
\renewcommand{\floatpagefraction}{0.75}

% --- Pandoc-Kompatibilität --------------------------------------------------
% Pandoc erzeugt diese Befehle in den Kapiteldateien.
\providecommand{\tightlist}{%
  \setlength{\itemsep}{0pt}\setlength{\parskip}{0pt}}
\makeatletter
\newsavebox\pandoc@box
\providecommand*\pandocbounded[1]{%
  \sbox\pandoc@box{#1}%
  \Gscale@div\@tempa{\textheight}{\dimexpr\ht\pandoc@box+\dp\pandoc@box\relax}%
  \Gscale@div\@tempb{\linewidth}{\wd\pandoc@box}%
  \ifdim\@tempb\p@<\@tempa\p@\let\@tempa\@tempb\fi
  \ifdim\@tempa\p@<\p@\scalebox{\@tempa}{\usebox\pandoc@box}%
  \else\usebox{\pandoc@box}%
  \fi}
\makeatother
\providecommand{\passthrough}[1]{#1}
\providecommand{\ul}[1]{\underline{#1}}
\providecommand{\st}[1]{\sout{#1}}
\usepackage[normalem]{ulem}
\usepackage{fancyvrb}
\usepackage{url}
\urlstyle{same}

% --- Literatur (biblatex + biber) ------------------------------------------
% Stil und Optionen setzt pdf.mjs aus .arbeit/einstellungen.md (Zitieren, stil).
% Nur geladen, wenn im Text zitiert wird: ohne Zitate braucht der Bau kein Biber.
\ifkitLiteratur
\makeatletter
\edef\kit@biblatex{\noexpand\usepackage[backend=biber,\kitBibOptionen]{biblatex}}
\kit@biblatex
\makeatother
\kitBibNachladen
\addbibresource{literatur.bib}
\fi

% --- Entwurfs-Wasserzeichen ------------------------------------------------
\ifkitEntwurf
  \DeclareNewLayer[
    background, textarea,
    contents={%
      \raisebox{-0.55\layerheight}[0pt][0pt]{%
        \makebox[\layerwidth]{%
          \rotatebox{50}{\fontsize{80}{80}\selectfont\color{black!7}\kitWortEntwurf}}}}
  ]{kit.entwurf}
  \AddLayersToPageStyle{@everystyle@}{kit.entwurf}
\fi
% Platzhalter für noch fehlende Unterkapitel im Entwurfsmodus.
\newcommand{\kitPlatzhalter}[1]{%
  \section{#1}%
  \begin{center}
    \fcolorbox{black!30}{black!3}{\parbox{0.85\linewidth}{\centering\small\itshape
      \kitWortPlatzhalter}}
  \end{center}}

% --- Links und Querverweise (immer zuletzt) --------------------------------
\usepackage[
  hidelinks,
  bookmarks=true,
  bookmarksnumbered=true,
  pdfstartview=FitH,
  pdftitle={\kitTitel},
  pdfauthor={\kitAutor}
]{hyperref}
\ifkitEnglisch
  \usepackage[english,capitalise,nameinlink]{cleveref}
\else
  \usepackage[ngerman,capitalise,nameinlink]{cleveref}
  \crefname{equation}{Gleichung}{Gleichungen}
  \Crefname{equation}{Gleichung}{Gleichungen}
\fi

% --- Eigene Ergänzungen der Nutzerin/des Nutzers ---------------------------
\IfFileExists{eigene-praeambel.tex}{\input{eigene-praeambel.tex}}{}


\begin{document}

% --- Titelei (römische Seitenzahlen) ---------------------------------------
\pagenumbering{Roman}
\ifkitTudscr
  % Deckblatt mit der TU-Dresden-Klasse tudscr (vorlage: tudscr unter Layout in .arbeit/einstellungen.md).
% Achtung (Stand 2026-09-30): tudscr setzt das alte Corporate Design um,
% das neue TU-Logo fehlt. Vor der Abgabe mit dem Lehrstuhl klären.
% Ungetestet im Kit, Klasse muss in MiKTeX/TeX Live installiert sein.
\faculty{\kitFakultaet}
\institute{\kitInstitut}
\title{\kitTitel}
\ifx\kitUntertitel\empty\else\subtitle{\kitUntertitel}\fi
\author{\kitAutor}
\ifx\kitMatrikel\empty\else\matriculationnumber{\kitMatrikel}\fi
\ifx\kitStudiengang\empty\else\course{\kitStudiengang}\fi
\ifx\kitBetreuer\empty\else\supervisor{\kitBetreuer}\fi
\ifx\kitErstgutachter\empty\else\referee{\kitErstgutachter}\fi
\ifx\kitZweitgutachter\empty\else\referee{\kitZweitgutachter}\fi
\date{\kitAbgabe}
\thesis{\kitTudscrThesis}
\maketitle

\else
  % Deckblatt. Alle Angaben kommen aus .arbeit/einstellungen.md (via kit-daten.tex).
% Zeilen ohne Wert blendet pdf.mjs automatisch aus.
\begin{titlepage}
  \thispagestyle{empty}
  \begin{center}
    \ifkitLogo
      \includegraphics[height=2.2cm,keepaspectratio]{\kitLogo}\par
      \vspace{0.8cm}
    \fi
    {\large \kitHochschule\par}
    \ifx\kitFakultaet\empty\else {\normalsize \kitFakultaet\par}\fi
    \ifx\kitInstitut\empty\else {\normalsize \kitInstitut\par}\fi

    \vspace{3.2cm}
    {\normalsize\MakeUppercase{\kitTyp}\par}
    \vspace{1.2cm}
    {\LARGE\bfseries \kitTitel\par}
    \ifx\kitUntertitel\empty\else
      \vspace{0.5cm}{\large \kitUntertitel\par}
    \fi

    \vfill
    \begin{tabular}{@{}rl@{}}
      \kitDeckblattZeilen
    \end{tabular}
    \vspace{1.5cm}
  \end{center}
\end{titlepage}

\fi
\ifkitSperrvermerk \chapter*{\kitWortSperrvermerk}
\thispagestyle{plain}
\ifkitEnglisch
This thesis contains confidential information. Publication, reproduction or
disclosure to third parties, in whole or in part, is not permitted without the
express written consent of the author and the cooperating institution. The
thesis is accessible only to the examiners and members of the examination board.
\else
Die vorliegende Arbeit enthält vertrauliche Informationen. Eine Veröffentlichung,
Vervielfältigung oder Weitergabe an Dritte, auch auszugsweise, ist ohne
ausdrückliche schriftliche Zustimmung der Verfasserin bzw. des Verfassers und der
beteiligten Einrichtung nicht gestattet. Die Arbeit ist ausschließlich den
Gutachtenden und den Mitgliedern des Prüfungsausschusses zugänglich.
\fi
\cleardoublepage
 \fi
\ifkitDanksagung % Text aus .arbeit/danksagung.md (von pdf.mjs nach danksagung-inhalt.tex gewandelt).
\chapter*{\kitWortDanksagung}
\thispagestyle{plain}
\input{danksagung-inhalt}
\cleardoublepage
 \fi
\ifkitZusammenfassung % Zusammenfassung und Abstract aus .arbeit/zusammenfassung.md
% (Abschnitte "## Zusammenfassung" und "## Abstract").
\ifkitHatZusammenfassung
  \chapter*{Zusammenfassung}
  \thispagestyle{plain}
  \begin{otherlanguage}{ngerman}
  \input{zusammenfassung-de}
  \end{otherlanguage}
  \cleardoublepage
\fi
\ifkitHatAbstract
  \chapter*{Abstract}
  \thispagestyle{plain}
  \begin{otherlanguage}{english}
  \input{zusammenfassung-en}
  \end{otherlanguage}
  \cleardoublepage
\fi
 \fi

\tableofcontents
\ifkitAbbildungen \listoffigures \fi
\ifkitTabellen \listoftables \fi
\ifkitAbkuerzungen % Abkürzungsverzeichnis. Einträge generiert pdf.mjs aus .arbeit/begriffe.md
% (alle Tabellenzeilen mit gefüllter Spalte "Abkürzung").
\chapter*{\kitWortAbkuerzungen}
\addcontentsline{toc}{chapter}{\kitWortAbkuerzungen}
\thispagestyle{plain}
\begin{longtable}{@{}p{3cm}p{\dimexpr\linewidth-3cm-2\tabcolsep\relax}@{}}
  \input{abkuerzungen-inhalt}
\end{longtable}
% longtable ohne Beschriftung zählt sonst den Tabellenzähler mit hoch.
\addtocounter{table}{-1}
 \fi
\ifkitFormelzeichen % Formelzeichen- und Symbolverzeichnis. Einträge generiert pdf.mjs aus
% .arbeit/begriffe.md (Tabelle mit Spalte "Symbol").
\chapter*{\kitWortFormelzeichen}
\addcontentsline{toc}{chapter}{\kitWortFormelzeichen}
\thispagestyle{plain}
\begin{longtable}{@{}p{2.5cm}p{\dimexpr\linewidth-5.5cm-4\tabcolsep\relax}p{3cm}@{}}
  \textbf{\kitWortSymbol} & \textbf{\kitWortBedeutung} & \textbf{\kitWortEinheit}\\
  \midrule
  \endhead
  \input{formelzeichen-inhalt}
\end{longtable}
% longtable ohne Beschriftung zählt sonst den Tabellenzähler mit hoch.
\addtocounter{table}{-1}
 \fi

% --- Textteil (arabische Seitenzahlen) -------------------------------------
\cleardoublepage
\pagenumbering{arabic}
\input{kapitel-liste}

% --- Literatur und Anhang --------------------------------------------------
\ifkitLiteratur
  \printbibliography[heading=bibintoc]
\fi
\ifkitAnhang
  \appendix
  \input{anhang-liste}
\fi
\ifkitHilfsmittel % Hilfsmittelverzeichnis und Erklärung zur Nutzung generativer KI.
% Aufbau nach dem in der Praxis bewährten KOGNITION-Schema: Verantwortung,
% verwendete Werkzeuge, Nutzung je Arbeitsschritt, Eigenleistung.
% Inhalte generiert pdf.mjs aus .arbeit/hilfsmittel.md.
\chapter*{\kitWortHilfsmittel}
\addcontentsline{toc}{chapter}{\kitWortHilfsmittel}
\thispagestyle{plain}

\section*{\kitWortVerantwortung}
{\small
\begin{tabularx}{\linewidth}{@{}X>{\centering\arraybackslash}p{1.2cm}@{}}
  \toprule
  \textbf{\kitWortAussage} & \textbf{\kitWortJa}\\
  \midrule
  \kitHmAussageEins & $\boxtimes$\\[4pt]
  \kitHmAussageZwei & $\boxtimes$\\[4pt]
  \kitHmAussageDrei & $\boxtimes$\\
  \bottomrule
\end{tabularx}}

\input{hilfsmittel-inhalt}
 \fi
\ifkitErklaerung % Selbstständigkeitserklärung. Der Wortlaut ist ein allgemeiner Standard.
% Viele Prüfungsordnungen schreiben einen eigenen Text vor: dann den Text
% unter .arbeit/latex/eigene-erklaerung.tex ablegen, er ersetzt diesen automatisch.
\IfFileExists{eigene-erklaerung.tex}{\input{eigene-erklaerung.tex}}{%
\chapter*{\kitWortErklaerung}
\addcontentsline{toc}{chapter}{\kitWortErklaerung}
\thispagestyle{plain}
\vspace{0.5cm}
\ifkitEnglisch
I hereby declare that I have written this \kitTyp{} independently and have used
no sources or aids other than those indicated. All passages taken verbatim or in
substance from published or unpublished works are marked as such. The use of
generative AI tools is fully documented in the list of aids. This thesis has not
been submitted in the same or a similar form to any other examination authority.
\else
Hiermit versichere ich, dass ich die vorliegende \kitTyp{} selbstständig verfasst
und keine anderen als die angegebenen Quellen und Hilfsmittel benutzt habe. Alle
Stellen, die wörtlich oder sinngemäß aus veröffentlichten oder nicht
veröffentlichten Schriften entnommen wurden, sind als solche kenntlich gemacht.
Die Nutzung generativer KI-Werkzeuge ist im Hilfsmittelverzeichnis vollständig
dokumentiert. Die Arbeit hat in gleicher oder ähnlicher Form noch keiner anderen
Prüfungsbehörde vorgelegen.
\fi

\vspace{2.5cm}
\noindent\kitOrt\ifx\kitOrt\empty\else, \fi\kitAbgabe

\vspace{2cm}
\noindent\rule{7cm}{0.4pt}\\
\kitAutor
}
 \fi

\end{document}
